\chapter{TINJAUAN PUSTAKA}
\label{pustaka}

\section{Isi Tinjauan Pustaka}
Ilmu pengetahuan merupakan produk budaya yang bersifat kumulatif, artinya ia merupakan karya dari banyak orang yang didokumentasikan dalam pustaka (misal buku teks, jurnal ilmiah, prosiding, laporan teknis/penelitian, majalah ilmiah dan dokumen paten). Oleh karena itu sebelum melakukan penelitian sebaiknya dilakukan studi terhadap pustaka yang terkait dengan tema yang akan diteliti untuk memperoleh \underline{\textbf{data/fakta}} tentang:

\begin{itemize}
    \item apa-apa yang sudah dilakukan oleh ilmuwan atau peneliti sebelumnya dengan sudut pandang atau aspek penelitian yang beragam beserta hasil-hasil yang diperolehnya, dan
\item apa-apa yang perlu diteliti lebih lanjut: (1) karena adanya pembatasan-pembatasan pada penelitian sebelumnya, atau (2) dengan sudut pandang atau aspek penelitian yang berbeda.
\end{itemize}



Dari hasil studi pustaka tersebut akan diperoleh gambaran mengenai langkah yang tepat untuk melaksanakan penelitian baik dari sisi sasaran/tujuan maupun metodologinya.

\underline{Perlu diperhatikan}, pustaka yang diacu harus dipastikan berasal dari sumber yang terpercaya. Untuk itu, peneliti harus bisa membedakan antara data/fakta dan opini/pendapat.

Hanya sumber yang memberikan informasi/fakta/data sajalah yang boleh diacu, sedangkan sumber yang hanya menyampaikan opini/pendapat tidak boleh diacu. Dengan demikian informasi yang diperoleh dari sumber manapun, termasuk internet, harus dipilah-pilah dan diambil hanya yang menyajikan data/fakta dengan benar didukung oleh penelitian, bukan sekedar opini/pendapat.

\section{Penulisan Acuan}
\textcolor{blue}{Sangat disarankan penulisan acuan (pustaka) dikelola menggunakan aplikasi manajemen pustaka, misalnya Zotero, Mendeley, atau lainnya.}

\textcolor{blue}{Kebaruan pustaka yang digunakan sebaiknya diperhatikan. Pustaka yang diacu sangat disarankan memiliki kebaruan, misalnya dengan memilih tahun penerbitan paling tua 20 tahun.}